\section*{Data and Code Availability}

Transparency and reproducibility are core principles of this work.
Data and code availability are described below,
subject to ethical, legal, and institutional constraints.

\subsection*{Data availability}

The data used in this study are [publicly available / available upon request / not publicly available].
If publicly available, specify the repository and license:
\textit{“The de-identified dataset used in this study is available at [repository link] under [license].”}

If data cannot be shared publicly, state the reason and access conditions:
\textit{“Patient-level data are protected under institutional and ethical approvals
and cannot be shared publicly.
Qualified researchers may request access through [data access committee or institution],
subject to appropriate approvals.”}

Where applicable, indicate whether aggregated or synthetic data are available.

\subsection*{Code availability}

Code used for data preprocessing, analysis, and model training is
[publicly available / partially available / not publicly available].

If available, provide a repository link and version reference:
\textit{“Code is available at \url{https://github.com/yourrepo/yourproject}
(commit \texttt{abc123}) for reproducibility.”}

If parts of the code cannot be shared,
briefly state which components are unavailable and why
(e.g., proprietary deployment interfaces),
and indicate whether pseudocode or detailed methodological descriptions are provided instead.

\section{Supplementary methods}

Include detailed methodological information that would disrupt the flow of the main text. Present additional analyses or derivations.

\section{Additional results}

Present supporting analyses and figures that complement your main results. Include any results referenced but not shown in the main text.
